\documentclass[../main.tex]{subfiles}
\begin{document}

\chapter{Хэрэгжүүлэлт}

Энэхүү бүлэгт системийн хэрэгжүүлэлтийн үндсэн хэсгүүд болох өгөгдлийн сангийн бүрдүүлэлт, загвар сургалт, серверийн хөгжүүлэлт, мобайл болон вэб интерфейсийн хэрэгжүүлэлтийг тайлбарлана.

\section{Өгөгдлийн сангийн бүрдүүлэлт}

Машин сургалтын загвар нь чанартай өгөгдлөөс шууд хамаардаг \cite{data_augmentation}. George болон Floerkemeier нарын GroZi-120 датасет \cite{grozi_dataset} нь жижиглэн худалдааны барааны анхны нээлттэй датасетуудын нэг байсан бөгөөд бид энэхүү судалгааны арга зүйг өөрийн датасет бүрдүүлэлтэд ашигласан.

\subsection{Өгөгдөл цуглуулах}

Өгөгдөл цуглуулах үйл явцыг дараах байдлаар зохион байгуулсан:

\begin{itemize}
\item \textbf{Зураг авах орчин:} Улаанбаатар хотын 5 өөр дэлгүүрийн лангуун
\item \textbf{Ашигласан төхөөрөмж:} iPhone 12 Pro (12MP), Samsung Galaxy S21 (64MP)
\item \textbf{Зургийн нягтрал:} 1920x1080 болон түүнээс дээш
\item \textbf{Гэрэлтүүлгийн нөхцөл:} Байгалийн гэрэл, LED гэрэлтүүлэг, холимог гэрэл
\item \textbf{Авах өнцөг:} Урд талаас, хажуугаас, дээрээс (30--60 градус)
\end{itemize}

\subsection{Өгөгдлийн статистик}

\begin{table}[H]
\centering
\caption{Өгөгдлийн сангийн статистик}
\label{tab:dataset_stats}
\begin{tabular}{|l|c|}
\hline
\textbf{Үзүүлэлт} & \textbf{Тоо} \\
\hline
Нийт зургийн тоо & 324 \\
\hline
Сургалтын багц (Train) & 227 (70\%) \\
\hline
Баталгаажуулалтын багц (Validation) & 65 (20\%) \\
\hline
Туршилтын багц (Test) & 32 (10\%) \\
\hline
Нийт ангиллын тоо & 8 \\
\hline
Нийт annotation тоо & 2,847 \\
\hline
Дундаж объект/зураг & 8.79 \\
\hline
\end{tabular}
\end{table}

\subsection{Барааны ангилал}

\begin{table}[H]
\centering
\caption{Барааны ангиллын жагсаалт}
\label{tab:product_classes}
\begin{tabular}{|c|l|c|c|}
\hline
\textbf{ID} & \textbf{Ангиллын нэр} & \textbf{Тоо} & \textbf{Хувь} \\
\hline
0 & Coca-Cola 500ml & 412 & 14.5\% \\
\hline
1 & Pepsi 500ml & 389 & 13.7\% \\
\hline
2 & Fanta 500ml & 356 & 12.5\% \\
\hline
3 & Sprite 500ml & 341 & 12.0\% \\
\hline
4 & Bon Aqua 500ml & 398 & 14.0\% \\
\hline
5 & Vittel 500ml & 287 & 10.1\% \\
\hline
6 & Rich Orange Juice 1L & 334 & 11.7\% \\
\hline
7 & Dobry Apple Juice 1L & 330 & 11.6\% \\
\hline
\multicolumn{2}{|c|}{\textbf{Нийт}} & \textbf{2,847} & \textbf{100\%} \\
\hline
\end{tabular}
\end{table}

\subsection{Шошгололт}

Цуглуулсан зургууд дээрх объектуудыг \texttt{LabelImg} хэрэгслийг ашиглан гар аргаар шошголж, YOLO форматаар хадгалсан. YOLO форматын жишээ:

\begin{lstlisting}[language=Python, caption=YOLO annotation формат]
# class_id x_center y_center width height
0 0.456 0.523 0.089 0.234
1 0.234 0.678 0.095 0.245
4 0.789 0.345 0.078 0.198
\end{lstlisting}

\subsection{Өгөгдөл баяжуулалт}

Сургалтын өгөгдлийг нэмэгдүүлэх зорилгоор дараах баяжуулалтын аргуудыг хэрэглэсэн:

\begin{table}[H]
\centering
\caption{Өгөгдөл баяжуулалтын аргууд}
\label{tab:augmentation}
\begin{tabular}{|l|l|c|}
\hline
\textbf{Арга} & \textbf{Тайлбар} & \textbf{Параметр} \\
\hline
Horizontal Flip & Хэвтээ тэнхлэгээр эргүүлэх & p=0.5 \\
\hline
Random Brightness & Гэрэлтүүлгийг өөрчлөх & $\pm$20\% \\
\hline
Random Rotation & Санамсаргүй эргүүлэх & $\pm$15 градус \\
\hline
Mosaic & 4 зургийг нэгтгэх & YOLOv8 default \\
\hline
MixUp & 2 зургийг холих & $\alpha$=0.5 \\
\hline
\end{tabular}
\end{table}

\section{Загвар сургалт}

\subsection{Сургалтын орчин}

\begin{table}[H]
\centering
\caption{Сургалтын орчны тохиргоо}
\label{tab:training_env}
\begin{tabular}{|l|l|}
\hline
\textbf{Бүрэлдэхүүн} & \textbf{Утга} \\
\hline
Үйлдлийн систем & Ubuntu 22.04 LTS \\
\hline
CPU & AMD Ryzen 7 5800X \\
\hline
GPU & NVIDIA RTX 3080 (10GB) \\
\hline
RAM & 32GB DDR4 \\
\hline
Python & 3.10.12 \\
\hline
PyTorch & 2.0.1+cu118 \\
\hline
Ultralytics & 8.0.200 \\
\hline
CUDA & 11.8 \\
\hline
\end{tabular}
\end{table}

\subsection{Гипер-параметрүүд}

\begin{table}[H]
\centering
\caption{Сургалтын гипер-параметрүүд}
\label{tab:hyperparameters}
\begin{tabular}{|l|c|l|}
\hline
\textbf{Параметр} & \textbf{Утга} & \textbf{Тайлбар} \\
\hline
Epochs & 100 & Давталтын тоо \\
\hline
Batch size & 16 & Багцын хэмжээ \\
\hline
Image size & 640$\times$640 & Оролтын хэмжээ \\
\hline
Learning rate (lr0) & 0.01 & Анхны суралцах хурд \\
\hline
Learning rate (lrf) & 0.001 & Эцсийн суралцах хурд \\
\hline
Momentum & 0.937 & SGD momentum \\
\hline
Weight decay & 0.0005 & L2 регуляризаци \\
\hline
IoU threshold & 0.7 & NMS IoU босго \\
\hline
\end{tabular}
\end{table}

\subsection{Сургалтын код}

\begin{lstlisting}[language=Python, caption=YOLOv8 загвар сургах код]
from ultralytics import YOLO

def train_model():
    model = YOLO("yolov8n.pt")
    results = model.train(
        data="models/configs/dataset.yaml",
        epochs=100,
        batch=16,
        imgsz=640,
        patience=20,
        save=True,
        device=0,
        project="runs/train",
        name="product_detector"
    )
    return results

if __name__ == "__main__":
    train_model()
\end{lstlisting}

\section{Серверийн хөгжүүлэлт}

Backend серверийг Python FastAPI ашиглан бүтээсэн. FastAPI нь асинхрон дэмждэг тул сүлжээний хүсэлтийг хурдан боловсруулдаг.

\subsection{Төслийн бүтэц}

\begin{lstlisting}[language=bash, caption=Backend-ийн файлын бүтэц]
backend/
|-- app/
|   |-- main.py              # FastAPI app
|   |-- config.py            # Тохиргоо
|   |-- database.py          # MongoDB холболт
|   |-- models/              # Pydantic схемүүд
|   |-- routes/              # API endpoints
|   |-- services/            # Бизнес логик
|-- requirements.txt
|-- Dockerfile
\end{lstlisting}

\subsection{Илрүүлэлтийн үйлчилгээ}

\begin{lstlisting}[language=Python, caption=ProductDetector класс]
from ultralytics import YOLO
import numpy as np
from typing import Dict

class ProductDetector:
    _instance = None

    def __new__(cls):
        if cls._instance is None:
            cls._instance = super().__new__(cls)
            cls._instance._initialize()
        return cls._instance

    def _initialize(self):
        self.model = YOLO("models/weights/best.pt")
        self.confidence_threshold = 0.25

    def detect(self, image: np.ndarray) -> Dict:
        results = self.model(image, conf=self.confidence_threshold)[0]
        detections = []
        for box in results.boxes:
            detections.append({
                "class_id": int(box.cls),
                "class_name": self.model.names[int(box.cls)],
                "confidence": float(box.conf),
                "bbox": box.xyxy[0].tolist()
            })
        return {"detections": detections, "total": len(detections)}
\end{lstlisting}

\section{Интерфэйсийн хэрэгжүүлэлт}

Системийн хэрэглэгчийн интерфейсийг хоёр үндсэн хэсэгт хуваан хэрэгжүүлсэн: мобайл аппликейшн болон вэб админ самбар.

\subsection{Мобайл аппликейшн}

Аудиторуудад зориулсан мобайл аппликейшнийг Odoo Mobile фреймворк дээр суурилан хэрэгжүүлсэн.

\subsubsection{Нэвтрэх хуудас}

\begin{figure}[H]
\centering
\includegraphics[width=0.32\textwidth]{images/application/Login_Page.png}
\caption{Мобайл аппликейшны нэвтрэх хуудас}
\label{fig:app_login}
\end{figure}

Нэвтрэх хуудас нь хэрэглэгчийн нэр, нууц үг оруулах талбар болон JWT баталгаажуулалтыг хэрэгжүүлсэн.

\subsubsection{Үндсэн нүүр хуудас}

\begin{figure}[H]
\centering
\includegraphics[width=0.32\textwidth]{images/application/Home_Page.png}
\caption{Мобайл аппликейшны нүүр хуудас}
\label{fig:app_home}
\end{figure}

Нүүр хуудас нь аудиторт өнөөдрийн хийх ажлуудын товч мэдээллийг харуулна.

\subsubsection{Кампанит ажлын хуудас}

\begin{figure}[H]
\centering
\includegraphics[width=0.32\textwidth]{images/application/Campaign_Page.png}
\caption{Кампанит ажлын жагсаалт}
\label{fig:app_campaign}
\end{figure}

\subsubsection{Зураг авах хуудас}

\begin{figure}[H]
\centering
\begin{subfigure}{0.45\textwidth}
\centering
\includegraphics[width=0.75\textwidth]{images/application/Image_Page.png}
\caption{Зураг авах интерфэйс}
\end{subfigure}
\hfill
\begin{subfigure}{0.45\textwidth}
\centering
\includegraphics[width=0.75\textwidth]{images/application/Image_Page_v2.png}
\caption{Зураг авсны дараах харагдац}
\end{subfigure}
\caption{Зураг авах функц}
\label{fig:app_image}
\end{figure}

\subsubsection{Дэлгүүрийн байршил}

\begin{figure}[H]
\centering
\includegraphics[width=0.32\textwidth]{images/application/Map_Page.png}
\caption{Дэлгүүрийн байршил}
\label{fig:app_map}
\end{figure}

Газрын зургийн хуудас нь Google Maps API ашиглан интеграци хийгдсэн.

\subsubsection{Аудитын түүх}

\begin{figure}[H]
\centering
\includegraphics[width=0.32\textwidth]{images/application/History_Page.png}
\caption{Аудитын түүхийн хуудас}
\label{fig:app_history}
\end{figure}

\subsection{Вэб админ самбар}

Удирдлагын ажилтнуудад зориулсан вэб интерфейсийг Odoo-ийн веб фреймворк дээр хэрэгжүүлсэн.

\subsubsection{Нүүр хуудас}

\begin{figure}[H]
\centering
\includegraphics[width=0.72\textwidth]{images/backoffice/Нүүр_хуудас.png}
\caption{Админ самбарын нүүр хуудас}
\label{fig:admin_dashboard}
\end{figure}

\subsubsection{Кампанит ажлууд}

\begin{figure}[H]
\centering
\includegraphics[width=0.72\textwidth]{images/backoffice/Кампанит ажлууд.png}
\caption{Кампанит ажлуудын жагсаалт}
\label{fig:admin_campaigns}
\end{figure}

\begin{figure}[H]
\centering
\includegraphics[width=0.72\textwidth]{images/backoffice/Кампанит_ажлууд_нэмэх.png}
\caption{Шинэ кампанит ажил үүсгэх}
\label{fig:admin_campaign_add}
\end{figure}

\subsubsection{Аудиторууд}

\begin{figure}[H]
\centering
\includegraphics[width=0.72\textwidth]{images/backoffice/Аудиторууд.png}
\caption{Аудиторуудын жагсаалт}
\label{fig:admin_auditors}
\end{figure}

\subsubsection{Дэлгүүрүүд}

\begin{figure}[H]
\centering
\includegraphics[width=0.72\textwidth]{images/backoffice/Дэлгүүрүүд.png}
\caption{Дэлгүүрүүдийн жагсаалт}
\label{fig:admin_stores}
\end{figure}

\subsubsection{Ангилалууд}

\begin{figure}[H]
\centering
\includegraphics[width=0.72\textwidth]{images/backoffice/Ангилалууд.png}
\caption{Барааны ангилалууд}
\label{fig:admin_categories}
\end{figure}

\subsubsection{Тайлан}

\begin{figure}[H]
\centering
\includegraphics[width=0.72\textwidth]{images/backoffice/Тайлан.png}
\caption{Аудитын тайлан}
\label{fig:admin_report}
\end{figure}

\begin{figure}[H]
\centering
\includegraphics[width=0.72\textwidth]{images/backoffice/Тайлан-1.png}
\caption{Дэлгэрэнгүй тайлан}
\label{fig:admin_report_detail}
\end{figure}

\section{Docker контейнержуулалт}

Системийг хялбар байршуулах зорилгоор Docker Compose ашиглан контейнержуулсан.

\begin{lstlisting}[language=bash, caption=docker-compose.yml]
version: '3.8'
services:
  mongodb:
    image: mongo:6.0
    ports:
      - "27017:27017"
    volumes:
      - mongodb_data:/data/db

  backend:
    build: ./backend
    ports:
      - "8000:8000"
    environment:
      - MONGODB_URI=mongodb://mongodb:27017
      - MODEL_PATH=/app/models/weights/best.pt
    depends_on:
      - mongodb

  frontend:
    build: ./frontend
    ports:
      - "3000:3000"
    depends_on:
      - backend

volumes:
  mongodb_data:
\end{lstlisting}

\end{document}
